\section{Обзор существующих решений}
\label{sec:Chapter2} \index{Chapter2}

%Здесь надо рассмотреть все существующие решения поставленной задачи, но не
%просто пересказать, в чем там дело, а оценить степень их соответствия тем
%ограничениям, которые были сформулированы в постановке задачи.

\subsection{Управление размещением (CPU- и NUMA-pinning)}

Это один из наиболее простых и широко используемых подходов, который заключается в явном управлении размещением процессов и памяти. Применяется в Linux с помощью утилиты \texttt{numactl}~\cite{numactl}, а также механизма \texttt{CPUSETS}~\cite{CPUSETS}. Утилита \texttt{numactl} позволяет задавать политику размещения вычислений и памяти, включая привязку процесса к определённым NUMA-нодам и управление распределением памяти между ними. Подсистема \texttt{CPUSETS} позволяет закреплять процессы за заданными наборами процессоров и нод памяти.

Данный подход обеспечивает высокий уровень контроля. При надлежащем применении в некоторых ситуациях можно сильно улучшить производительность приложений, в том числе уменьшив количество обращений к удалённой памяти. Однако есть и недостатки:
\begin{enumerate}
    \item\label{itm:pinning1} Иногда, чтобы действительно добиться улучшения производительности, требуется понимание архитектуры и работы приложения. Для этого часто необходим анализ, например, через профилирование. После этого механизм требуется ещё вручную настроить. Все это может быть не очень удобно для конечного пользователя.
    \item\label{itm:pinning2} Привязка статична~--- она ограничивает гибкость планировщика и не адаптируется к изменениям поведения приложения и состояния системы со временем.
    \item\label{itm:pinning3} Политики размещения памяти помогают сделать её выделение более равномерным и/или предсказуемым, но все равно не решают проблему удаленного доступа из-за несоответствия между потоками и используемыми ими данными.
\end{enumerate}

Решение, компенсирующее недостатки~\ref{itm:pinning1} и ~\ref{itm:pinning2}, уже существует~--- это демон \texttt{numad}~\cite{numad}. Заключается оно в том, что \texttt{numad}-демон в реальном времени отслеживает процессы в системе, анализирует память и CPU, используемые ими. И, исходя из этого, принимает решение о размещении наиболее долгоживущих и ресурсоемких процессов в NUMA-системе для повышения эффективности использования памяти и CPU. Т.е. получается такой автоматизированный CPU- и NUMA-pinning.

Однако основная проблема заключается в другом: CPU- и NUMA-pinning, как ручной, так и автоматический, в принципе плохо подходят для больших cross-NUMA приложений, оптимизация которых, собственно, и является целью работы. Дело в том, что cross-NUMA приложения, они на то и cross-NUMA, что не могут уместиться на одной NUMA-ноде. В общем случае они будут использовать все NUMA-ноды, и привязка потеряет какой-либо смысл. Она в большей степени предназначена для серверов, где может быть запущено несколько приложений, виртуальных машин, например~--- и их уже можно пытаться группировать по NUMA-нодам, добиваясь баланса.

Остается лишь пробовать применять различные политики размещения памяти, но проблема удаленного доступа, описанная в недостатке~\ref{itm:pinning3}, продолжает быть актуальной.

\subsection{Автоматическая NUMA-балансировка в ядре Linux}
\label{sec:Chapter2_NUMA_balancer}

Современные версии ядра Linux используют механизм автоматической NUMA-балансировки~\cite{numa_balancing_lwn, numa_balancing_doc}, включённый по умолчанию. Он концептуально работает следующим образом:
\begin{enumerate}
    \item периодически сканирует страницы процесса, инвалидируя соответствующие записи PTE-таблиц~--- это приводит к \textit{NUMA hinting faults} при последующих обращениях;
    \item при обработке \textit{fault}-а фиксирует NUMA-ноду, на которой произошёл доступ, собирает статистику, какие страницы где используются, насколько часто, какие CPU к ним обращаются;
    \item на основе статистики принимает решение о миграции страниц между NUMA-нодами, руководствуясь соображениями локальности.
\end{enumerate}

Этот подход является полностью прозрачным для пользовательских приложений и не требует их модификации. Он динамически адаптируется к поведению нагрузки. Тем не менее, балансировщик может быть недостаточно оптимален для больших NUMA-нагрузок:
\begin{enumerate}
    \item\label{itm:balance1} если несколько потоков на разных NUMA-нодах активно используют одни и те же данные, то миграция данных не даст выгоды, только накладные расходы;
    \item балансировщик преимущественно работает только с \textit{анонимной} памятью процесса, \textit{не разделяемой} с другими процессами. Он игнорирует \textit{read-only file-backed} страницы, т.к. это могут быть \textit{разделяемые} библиотеки, используемые всеми процессами в системе, и их миграция приведет к недостатку~\ref{itm:balance1}.
\end{enumerate}

Предлагаемый в данной работе механизм репликации таблиц трансляций и данных предназначен, чтобы разрешить недостатки NUMA-балансировщика:
\begin{enumerate}
    \item репликация вместо миграции~--- у каждой NUMA-ноды будет своя локальная копия;
    \item механизм репликации будет работать в том числе и с \textit{file-backed} памятью, \textit{разделяемой} несколькими процессами~--- об этом будет подробнее сказано в следующих главах.
\end{enumerate}

\subsection{Self-Replicating Page-Tables}

Стоит отметить, что исследования NUMA-эффекта, вызываемого таблицами трансляций, уже велись, и был предложен механизм репликации таблиц трансляций~\cite{Mitosis}. В дальнейшем идея получила развитие, и спустя несколько лет вышел механизм автоматического включения/выключения репликации, в зависимости от характеристик приложения~\cite{WASP}.

Таким образом, данная работа является следующим этапом развития идеи репликации памяти, добавляя возможность репликации не только таблиц трансляций, но и данных приложений.
